\chapter{Registry, Persistence, and Transactional State Management}
\label{chap:registry-persistence}

The GlobalPlatform model described in Chapter~4 assumes that a secure
element can remember what has been loaded, installed, personalized, keyed
and made selectable. Oxide SE makes this assumption explicit through one
kernel registry and one persistence strategy designed for small embedded
targets.

This chapter describes the architectural contract rather than a particular
board implementation. Chapter~8 gives the implementation point of view;
this chapter defines what the implementation must preserve.

%--------------------------------------------------
\section{Persistent Memory Model}
\label{sec:persistence-memory-model}

Oxide SE targets secure elements whose non-volatile storage is usually
flash-like memory. Such memory is persistent, may also be executable, and
is not byte-addressable in the same sense as RAM. The architecture
therefore assumes the following properties.

\begin{itemize}
\item The firmware image may execute directly from flash, for example in an
  execute-in-place configuration.
\item A board may reserve a distinct flash area for mutable registry
  persistence. This area must not overlap the immutable firmware image.
\item Flash is erased by pages or sectors. After erase, bytes read as
  \texttt{0xFF}.
\item Programming can write a page or a subrange, but it must be treated as
  non-atomic at the registry level.
\item Power may be lost at any point during an erase or write operation.
\end{itemize}

The failure model is deliberately conservative. After a reset, a block that
was being written when power disappeared may be partially programmed, may
contain a mixture of old and new bytes, or may look erased in places. No
registry state is considered valid merely because its header is present.
Every committed persistent block must therefore carry enough information to
validate itself independently.

The fundamental validity rule is:

\begin{quote}
A persistent block is usable only if its magic value is known, its advertised
length is page-aligned and inside the reserved persistence area, and its final
checksum validates the complete block.
\end{quote}

This rule is intentionally simple. It avoids relying on hidden hardware
transaction semantics and makes the same recovery logic applicable to real
flash, emulated flash and future EEPROM-like backends.

%--------------------------------------------------
\section{Logical Registry Contract}
\label{sec:registry-logical-contract}

The registry is the authoritative catalogue of managed objects inside the
secure element. It is a logical GlobalPlatform registry, not a collection of
independent per-application stores. Security Domains define authority and
visibility over registry objects, but the catalogue itself is one coherent
kernel structure.

Each registry entry has a common identity header:

\begin{itemize}
\item \textbf{object AID}: the stable identifier of the managed object;
\item \textbf{parent Security Domain AID}: the administrative Security
  Domain under whose authority the object exists;
\item \textbf{kind}: the object family, plus kind-specific state bits;
\item \textbf{data reference}: a reference to the persistent payload, when
  such a payload exists.
\end{itemize}

The object kind determines how the payload is interpreted. Oxide SE uses
the following managed-object families.

\begin{center}
\begin{tabular}{|l|p{8.5cm}|}
\hline
\textbf{Kind} & \textbf{Persistent meaning} \\
\hline
Security Domain &
Administrative instance with a package binding, base privilege bytes and
serialized Security-Domain state. \\
\hline
Package &
Executable Rustlet package, normally a FAE image. A package may be embedded
in immutable firmware or later stored in mutable persistence. \\
\hline
Instance &
Installed Rustlet instance with a package binding and serialized application
state. \\
\hline
Key &
Cryptographic key object owned by a Security Domain, including key metadata
such as type, version, identifier, usage and raw key material. \\
\hline
Data &
Generic Security-Domain-scoped data object, typically written through
\texttt{STORE DATA} and returned through \texttt{GET DATA} when no
standard or Security-Domain-specific tag handler claims the request. \\
\hline
\end{tabular}
\end{center}

This representation separates three concerns that are easy to confuse:

\begin{itemize}
\item the AID identifies the object;
\item the parent Security Domain defines the administrative authority;
\item the payload contains the state or bytes specific to the object kind.
\end{itemize}

Consequently, a Rustlet instance does not itself carry Security-Domain
privileges. Privilege bytes are properties of Security Domain objects.
Ordinary Rustlet instances carry only their package binding and serialized
state. Key objects are associated with a Security Domain through the same
parent-AID relation as other managed objects.

\subsection{GET DATA, STORE DATA, and Registry Fallback}
\label{subsec:registry-get-store-data-fallback}

\texttt{GET DATA} and \texttt{STORE DATA} are registry-facing management
commands, but not every \texttt{GET DATA} tag is treated as a generic data
object. Some tags are part of the Security Domain administrative interface
itself. For example, the card recognition data tag \texttt{0066} and the
card capability information tag \texttt{0067} are interpreted by the active
Security Domain backend, while the life-cycle tag \texttt{9F70} is derived
from its managed state. A stored registry data object must not override these
special tags. Section~\ref{sec:gp-get-data-discovery} defines the public GP
encoding of \texttt{0066} and \texttt{0067}.

For all other tags, Oxide SE applies a conservative fallback rule. The
active authority is resolved first: clear management commands are interpreted
under the root Security Domain authority, while protected commands are
interpreted under the Security Domain bound to the authenticated secure
channel. \texttt{STORE DATA} then writes a generic \textbf{Data} object owned
by that authority. A later \texttt{GET DATA} for the same non-special tag may
return that object if the active Security Domain does not provide a more
specific answer.

Rustlet-backed Security Domains follow the same rule with one additional
delegation step. The kernel first asks the selected Rustlet Security Domain
whether it handles the requested tag. If the Rustlet Security Domain returns
``unsupported'' for a non-special tag, the kernel may fall back to the
registry-backed \textbf{Data} object. If the tag is special, no registry
fallback is applied: the Security Domain answer remains authoritative.

\texttt{GET STATUS} is a distinct registry projection rather than a generic
data-object lookup. It enumerates packages, ordinary instances and Security
Domains as GP \texttt{E3} records, subject to the active authority's hierarchy
and privileges. Internal \textbf{Key} and \textbf{Data} objects are deliberately
excluded: discovery may report cryptographic capabilities, but it never
discloses key material or private registry values. Continuation state records
only the next stable registry slot and the canonical query; it does not retain
a copied registry object or response page.

Persistence is the storage mechanism behind this logical registry. It does
not change the registry semantics. A registry mutation such as installing an
instance, updating serialized state, loading a package, adding a key or
deleting an object first changes the logical registry state. The persistence
layer then publishes a durable representation of that new state.

%--------------------------------------------------
\section{Append-Only Transactional Registry}
\label{sec:append-only-transactional-registry}

Oxide SE uses an append-only flash journal for the mutable registry. The
goal is not to implement a complex filesystem. The goal is to make registry
updates robust against the one failure that matters most for small secure
elements: power loss during a flash operation.

\subsection{Personalization Marker}
\label{subsec:persistence-personalization-marker}

The first page of the reserved persistence area carries a personalization
marker. If this marker is absent or invalid, the secure element treats the
registry as unpersonalized. It then replays the predeployment manifest,
builds the initial registry in RAM, publishes that registry to flash and
finally leaves the marker in place.

The marker is not a registry block. It answers only one question: has the
mutable persistence area already been initialized for this image? The
registry state itself is always recovered from the latest valid registry
block.

\subsection{Block Families}
\label{subsec:persistence-block-families}

Every persistent object is stored as a page-aligned block. Each block starts
with a magic word and ends with a checksum. The currently defined block
families are:

\begin{center}
\begin{tabular}{|l|l|p{7.5cm}|}
\hline
\textbf{Magic} & \textbf{Name} & \textbf{Purpose} \\
\hline
\texttt{0x600DB055} & \texttt{BOSS} &
Registry block. The valid block with the highest mutation counter is the
authoritative registry. \\
\hline
\texttt{0x600DC0DE} & \texttt{C0DE} &
Executable package payload, typically a FAE image. \\
\hline
\texttt{0x600DBA5E} & \texttt{BA5E} &
Generic registry data, including key-like binary payloads when appropriate. \\
\hline
\texttt{0x600DFACE} & \texttt{FACE} &
Serialized runtime state, including Rustlet instance state and
Security-Domain state. \\
\hline
\end{tabular}
\end{center}

The registry block has the following logical layout:

\begin{center}
\begin{tabular}{|l|p{8.5cm}|}
\hline
\textbf{Field} & \textbf{Meaning} \\
\hline
magic & \texttt{0x600DB055}. \\
\hline
total length & Page-aligned length of the complete block. \\
\hline
mutation counter & Monotonic counter incremented for each committed
registry mutation. \\
\hline
entry count & Number of registry entries stored in this block. \\
\hline
entries & Fixed-size registry entries. Each entry contains encoded object
AID, encoded parent Security Domain AID, kind word and data reference. \\
\hline
checksum & Final checksum over the complete block except the checksum field
itself. \\
\hline
\end{tabular}
\end{center}

Object blocks use the simpler layout:

\begin{center}
\begin{tabular}{|l|p{8.5cm}|}
\hline
\textbf{Field} & \textbf{Meaning} \\
\hline
magic & One of \texttt{C0DE}, \texttt{BA5E} or \texttt{FACE}. \\
\hline
total length & Page-aligned length of the complete block. \\
\hline
payload length & Exact useful length of the payload. \\
\hline
payload & Object bytes interpreted according to the block family and
registry kind. \\
\hline
checksum & Final checksum over the complete block except the checksum field
itself. \\
\hline
\end{tabular}
\end{center}

\subsection{Dynamic Package Loading}
\label{subsec:persistence-dynamic-package-loading}

Dynamic executable loading is represented in the persistent registry by
\texttt{C0DE} object blocks. The GlobalPlatform command sequence is described
from the APDU point of view in Chapter~\ref{ch:communication-protocol}; this
section describes only the registry-side effect of that sequence.

\texttt{INSTALL [for load]} opens a volatile load transaction. It announces
the package AID, the expected package size and the expected package hash. If
the selected Security Domain authorizes the load, the kernel reserves one
page-aligned \texttt{C0DE} span large enough to hold the full package payload
and its persistent block metadata. This reservation is not a registry commit:
no \texttt{BOSS} block references the new package yet.

Each following \texttt{LOAD} command appends one consecutive fragment into the
reserved \texttt{C0DE} payload area. The persistent object header remains
kernel metadata. The executable FAE image starts after that header, at the
common object-payload offset, and the Rustlet runtime is later given exactly
that payload slice. The kernel therefore never branches to the \texttt{C0DE}
magic word or to the persistent length fields.

The final \texttt{LOAD} command may publish the package only if all load
invariants hold:

\begin{itemize}
\item block numbers were received consecutively;
\item the total number of received bytes is exactly the size announced by
  \texttt{INSTALL [for load]};
\item no fragment wrote beyond the reserved package payload;
\item the computed package hash matches the announced hash;
\item the final \texttt{C0DE} checksum validates.
\end{itemize}

Only after those checks does the kernel insert the Package object into the
in-RAM registry and publish a new \texttt{BOSS} block that references the
\texttt{C0DE} block. An interrupted load, a wrong hash, an incomplete final
block or an invalid checksum leaves no committed Package object behind. Such a
\texttt{C0DE} span is simply an unreferenced object block and can be recycled
by the normal journal rules.

Predeployed packages follow the same logical registry model, but their
payloads may live in the immutable firmware image rather than in the mutable
registry flash area. Registry references therefore distinguish firmware-image
payloads from registry-flash payloads while keeping the same logical Package
object kind.

\subsection{Commit Rule}
\label{subsec:persistence-commit-rule}

The transaction rule is intentionally strict:

\begin{quote}
The payload block for a new or modified object is written first. Only after
that payload is durable may a new \texttt{BOSS} block be written to reference
it.
\end{quote}

The existing \texttt{BOSS} block is never overwritten in place. If power is
lost while the new payload is being written, the old \texttt{BOSS} still
points to the old object. If power is lost while the new \texttt{BOSS} is
being written, its checksum will not validate and recovery will again select
the previous valid \texttt{BOSS}.

At boot, recovery scans the complete reserved persistence area. Erased pages
are valid free pages, not end markers. The scanner therefore keeps walking
after erased pages, because older valid blocks may still exist later in the
circular journal. Among all valid registry blocks found, the block with the
highest mutation counter is selected as the authoritative registry.

\subsection{Recycling Rule}
\label{subsec:persistence-recycling-rule}

Old blocks may be recycled only when they are no longer needed by the latest
valid registry. In the common case, this means:

\begin{itemize}
\item the latest valid \texttt{BOSS} block is protected;
\item every object block referenced by that \texttt{BOSS} is protected;
\item erased pages may be reused;
\item older valid blocks not referenced by the latest \texttt{BOSS} may be
  reused if their whole page-aligned span is large enough.
\end{itemize}

There is one additional transactional invariant. A registry update is one
logical mutation followed by one new \texttt{BOSS} publication. During that
publication, the previous \texttt{BOSS} remains the recovery root. Therefore
the old payload replaced or deleted by that one mutation must not be
considered free until the new \texttt{BOSS} has been completely written and
validated. Since Oxide SE does not batch several registry mutations before a
registry publication, there can be at most one such old-only payload per
in-flight transaction.

This rule is the small but important distinction between ``not referenced by
the new registry'' and ``safe to erase now''. A block becomes safely
recyclable only after there exists a newer valid \texttt{BOSS} whose recovery
does not require that block.
